\section{Method}
\label{sec:method}

\subsection{Research Goal}

This study compares four regression models---MLP, Decision Tree (DT), Random Forest
(RF), and mGBDT---on the \texttt{price} target, and asks whether richer feature
encoding or model structure improves over a naive baseline (a Random Forest on
ordinal-encoded categories).
Two axes are crossed under identical preprocessing and evaluation: (i) the
\emph{location-encoding variant} (\texttt{cat}, \texttt{tgt}, \texttt{coord\_only},
\texttt{tgt\_only}), and (ii) whether each model is trained as a single price
predictor or as a price / price-per-sqft \emph{stacking ensemble}. This yields
$4\ \text{models}\times 4\ \text{variants}\times 2\ \text{modes}=32$ configurations,
each evaluated by 5-fold cross-validation.

% -------------------------------------------------------
\subsection{Preprocessing}
% -------------------------------------------------------

Raw features are transformed as follows; all fold-dependent steps are fit on the
training rows only.

\textbf{Outliers and continuous features.}
Records with non-positive price and the top $0.1\%$ of prices are dropped,
leaving 4,086 records. \texttt{sqft\_lot} is replaced by
$\texttt{log\_sqft\_lot}=\log(1+\texttt{sqft\_lot})$ to reduce skew, and the
coordinates \texttt{lat}/\texttt{long} are projected onto the unit sphere
(\emph{coordinate encoding}):
\begin{equation}
  x = \cos(\text{lat})\cos(\text{lon}),\quad
  y = \cos(\text{lat})\sin(\text{lon}),\quad
  z = \sin(\text{lat}),
\end{equation}
which avoids the $\pm180°$ meridian discontinuity; missing coordinates use the
per-city (else global) mean.

\textbf{Temporal and scaling.}
We derive $\texttt{age}=\texttt{sale\_year}-\texttt{yr\_built}$ and a binary
$\texttt{was\_renovated}$, then drop the raw year columns. For scaling,
\texttt{x}, \texttt{y}, \texttt{z}, \texttt{condition}, \texttt{age},
\texttt{bedrooms}, \texttt{bathrooms}, \texttt{floors}, \texttt{view} are
MinMax-scaled to $[0,1]$, while \texttt{sqft\_living}, \texttt{sqft\_above},
\texttt{sqft\_basement}, \texttt{log\_sqft\_lot} and the encoded \texttt{city},
\texttt{zipcode} are Z-scored.

\textbf{Categorical and target.}
\texttt{city} and \texttt{zipcode} are encoded either as ordinal codes (naive
baseline) or via leakage-free \emph{target encoding}---each category replaced by
a smoothed mean of training price (smoothing $10$). \texttt{street},
\texttt{statezip}, \texttt{country}, \texttt{date} are removed as redundant. The
\texttt{price} target is MinMax-scaled for training and inverted to USD before
scoring (no log-transform), so all errors stay interpretable in dollars; the
auxiliary $\texttt{price\_per\_sqft}=\texttt{price}/\texttt{sqft\_living}$ is used
only by the stacking ensemble (Section~\ref{subsec:ensemble}).

% -------------------------------------------------------
\subsection{Location-Encoding Variants}
\label{subsec:variants}
% -------------------------------------------------------

We isolate the contribution of each location signal with four feature variants,
all sharing the same split and scaling:
\begin{itemize}
  \item \texttt{cat} --- ordinal codes for \texttt{city}/\texttt{zipcode}, no
    coordinates (the naive baseline encoding).
  \item \texttt{tgt} --- target-encoded \texttt{city}/\texttt{zipcode} \emph{plus}
    the cartesian coordinates \texttt{x,y,z}.
  \item \texttt{coord\_only} --- coordinates \texttt{x,y,z} only
    (\texttt{tgt} with \texttt{city}/\texttt{zipcode} dropped).
  \item \texttt{tgt\_only} --- target-encoded address only
    (\texttt{tgt} with \texttt{x,y,z} dropped).
\end{itemize}

% -------------------------------------------------------
\subsection{Model Configuration}
% -------------------------------------------------------

All four models use the fixed configurations in Table~\ref{tab:hp_grid}, which
produce every reported number. (Exploratory tuning scripts exist under
\texttt{experiments/}, but the comparison holds configurations fixed so that
differences are attributable to encoding and stacking, not per-cell tuning.)
The \textbf{DT} is a single CART regressor~\cite{breiman1984cart}
(interpretable lower bound); the \textbf{RF} is a bagged ensemble of 100 unpruned
trees, and RF/\texttt{cat} is the headline baseline; the \textbf{MLP} is a
$[32,32]$ ReLU network trained with Adam. The \textbf{mGBDT} uses XGBoost
forward/inverse modules~\cite{chen2016xgboost} under target propagation; the
reported \emph{single}-layer configuration reduces to standard XGBoost (a
two-layer variant performed worse, Section~\ref{sec:result}), so it is labelled
``mGBDT (XGBoost)''.

\begin{table}[h]
\centering
\caption{Fixed model configurations used for all reported runs.}
\label{tab:hp_grid}
\begin{tabular}{llll}
\toprule
\textbf{Model} & \multicolumn{3}{l}{\textbf{Configuration}} \\
\midrule
DT    & \texttt{max\_depth}$=6$ & \texttt{min\_samples\_split}$=2$ & \texttt{min\_samples\_leaf}$=2$ \\
\midrule
RF    & \texttt{n\_estimators}$=100$ & \texttt{max\_depth}$=$None & \texttt{min\_samples\_leaf}$=2$ \\
\midrule
MLP   & hidden $[32,32]$, Adam & lr $=10^{-2}$, batch $16$ & $\le 50$ epochs, patience $5$ \\
\midrule
mGBDT & 1 TP layer over XGBoost & lr $=0.1$, \texttt{max\_depth}$=3$ & 5 boost rounds, $20$ epochs \\
\bottomrule
\end{tabular}
\end{table}

% -------------------------------------------------------
\subsection{Price / Price-per-sqft Stacking Ensemble}
\label{subsec:ensemble}
% -------------------------------------------------------

To dampen the dominance of size and the influence of price outliers, each model can
also be trained as a two-member stacking ensemble. On every fold we train one model
to predict \texttt{price} and a second, identical model to predict
\texttt{price\_per\_sqft}; the second model's prediction is converted back to a price
estimate $\hat{p}_{2}=\widehat{\texttt{pps}}\times\texttt{sqft\_living}$. The final
prediction blends the two,
\begin{equation}
  \hat{p} = w\,\hat{p}_{1} + (1-w)\,\hat{p}_{2},
\end{equation}
where the weight $w\in[0,1]$ is grid-searched (101 steps) to minimize RMSE on that
fold's \emph{training} split before being applied to its validation split.

% -------------------------------------------------------
\subsection{Experimental Protocol}
% -------------------------------------------------------

All results use 5-fold cross-validation (shuffled, seed $42$) over the 4,086
cleaned records, with fold indices shared across every model, variant, and mode
for fairness. Target encoding, scaling, and the ensemble blend weight are fit on
training rows only, so no validation information leaks. For each configuration we
report the across-fold \emph{mean $\pm$ std} of every metric on both splits; the
single \texttt{price} target is used throughout.

% -------------------------------------------------------
\subsection{Evaluation Metrics}
% -------------------------------------------------------

We report five complementary regression metrics. RMSE squares residuals and so
penalises large mis-pricings; MAE weights all residuals equally and is
outlier-robust (a large RMSE--MAE gap flags outlier dominance); MAPE is a
scale-free percentage error; $R^2$ is the explained variance fraction; and
Adjusted $R^2$ corrects $R^2$ for the number of predictors $k$ for fair
comparison across models:
\begin{equation}
  \text{RMSE} = \sqrt{\tfrac{1}{n}\textstyle\sum_i(y_i - \hat{y}_i)^2},\quad
  \text{MAE} = \tfrac{1}{n}\textstyle\sum_i|y_i - \hat{y}_i|,\quad
  \text{MAPE} = \tfrac{1}{n}\textstyle\sum_i\big|\tfrac{y_i - \hat{y}_i}{y_i}\big|\times 100\%,
\end{equation}
\begin{equation}
  R^2 = 1 - \frac{\sum_i(y_i-\hat{y}_i)^2}{\sum_i(y_i-\bar{y})^2},\qquad
  \bar{R}^2 = 1 - \frac{(1-R^2)(n-1)}{n-k-1}.
\end{equation}
We rank configurations primarily by \textbf{Adjusted R$^2$ on the validation split}.
Secondary efficiency indicators (training time, inference time, peak memory) are not
recorded in the current pipeline and are left to future work.
